\section{Neural LoFi Kernel}

\label{2605:app:kernel-nlofi}

\subsection{Preliminaries: RKHS, spectral theorem and kernel integral operators}\label{2605:app:rkhs-prelims}

This subsection collects the standard background on Reproducing Kernel Hilbert Spaces (RKHS) and the spectral decomposition of kernel integral operators that we use in the rest of the appendix.
We refer to \cite{scholkopf2002learning,steinwart2008support} for a more detailed exposition. Throughout, $(\mathcal{X},\mu)$ denotes a measurable input space equipped with a finite Borel measure $\mu$ (e.g.\ the data distribution), and $K:\mathcal{X}\times\mathcal{X}\to\mathbb{R}$ a symmetric kernel.

\paragraph{Reproducing Kernel Hilbert Space}
A Hilbert space $\mathcal{H}$ of real-valued functions on $\mathcal{X}$ is a \emph{Reproducing Kernel Hilbert Space} (RKHS) if, for every $\bx\in\mathcal{X}$, the evaluation functional $\mathrm{ev}_\bx:f\mapsto f(\bx)$ is bounded (continuous) on $\mathcal{H}$. By the Riesz representation theorem, there exists a unique element $K_\bx\in\mathcal{H}$ such that
\begin{equation}\label{2605:eq:reproducing}
  f(\bx)=\langle f,K_\bx\rangle_{\mathcal{H}}\qquad\text{for all }f\in\mathcal{H}.
  \end{equation}
  Defining $K(\bx,\bx'):=\langle K_\bx,K_{\bx'}\rangle_{\mathcal{H}}=K_{\bx'}(\bx)$ yields a symmetric positive semi-definite kernel, and \eqref{2605:eq:reproducing} is called the \emph{reproducing property} because the kernel section $K_\bx=K(\bx,\cdot)$ literally reproduces the value of $f$ at $\bx$ via an inner product. Conversely, by the Moore--Aronszajn theorem~\cite{aronszajn1950theory}, every symmetric positive semi-definite kernel $K$ uniquely determines an RKHS $\mathcal{H}_K$ in which \eqref{2605:eq:reproducing} holds; concretely, $\mathcal{H}_K$ is obtained by completing $\mathrm{span}\{K(\bx,\cdot):\bx\in\mathcal{X}\}$ under the inner product $\langle K(\bx,\cdot),K(\bx',\cdot)\rangle_{\mathcal{H}_K}=K(\bx,\bx')$.

  \paragraph{Spectral theorem for compact self-adjoint operators}
  Let $\mathcal{H}$ be a separable Hilbert space and $T:\mathcal{H}\to\mathcal{H}$ a bounded linear operator. Recall that $T$ is \emph{self-adjoint} if $\langle Tf,g\rangle=\langle f,Tg\rangle$, and \emph{compact} if it maps bounded sets to relatively compact sets (equivalently, it is a norm-limit of finite-rank operators). The spectral theorem states:
  \begin{theorem}[Hilbert--Schmidt spectral theorem]\label{2605:thm:spectral}
  If $T$ is compact and self-adjoint on a separable Hilbert space $\mathcal{H}$, then there exist a (finite or countable) sequence of real eigenvalues $\{\lambda_i\}_{i\geq 1}$ with $|\lambda_1|\geq|\lambda_2|\geq\cdots\to 0$ and an orthonormal system $\{e_i\}_{i\geq 1}\subset\mathcal{H}$ of eigenvectors $Te_i=\lambda_i e_i$ such that
  \begin{equation}\label{2605:eq:spectral}
    T f  =  \sum_{i\geq 1} \lambda_i \langle f,e_i\rangle e_i\qquad\text{for all }f\in\mathcal{H}.
    \end{equation}
    \end{theorem}
  Trace-class and Hilbert--Schmidt operators are special cases of compact operators on which the spectrum is, respectively, absolutely summable ($\sum_i|\lambda_i|<\infty$) and square-summable ($\sum_i\lambda_i^2<\infty$).

    \paragraph{Kernel integral operator}
    Given a kernel $K$ and the measure $\mu$, the associated \emph{integral operator} $T_K:L^2(\mathcal{X},\mu)\to L^2(\mathcal{X},\mu)$ is defined by
    \begin{equation}\label{2605:eq:integral-op}
      (T_K f)(\bx) := \int_{\mathcal{X}} K(\bx,\bx') f(\bx') d\mu(\bx').
      \end{equation}
      Under the standard assumptions that $K$ is symmetric, continuous (or merely measurable) and square-integrable in the sense that $\int_{\mathcal{X}\times\mathcal{X}} K(\bx,\bx')^2 d\mu(\bx)d\mu(\bx')<\infty$, the operator $T_K$ is self-adjoint and Hilbert--Schmidt (in particular compact) on $L^2(\mathcal{X},\mu)$.

      Applying Theorem~\ref{2605:thm:spectral} to $T_K$ yields eigenpairs $\{(\lambda_i,e_i)\}_{i\geq 1}$ with $\lambda_i\geq 0$, $\lambda_i\downarrow 0$, and $\{e_i\}$ orthonormal in $L^2(\mathcal{X},\mu)$. Mercer's theorem~\cite{mercer1909functions} then states that, under mild continuity assumptions on $\mathcal{X}$ and $K$ (e.g.\ $\mathcal{X}$ compact and $K$ continuous), the kernel itself admits the absolutely and uniformly convergent expansion
      \begin{equation}\label{2605:eq:mercer}
        K(\bx,\bx') = \sum_{i\geq 1}\lambda_i e_i(\bx) e_i(\bx').
        \end{equation}

        \paragraph{Relation between the integral operator and the RKHS}
        The Mercer decomposition~\eqref{2605:eq:mercer} provides an explicit isometric description of the RKHS $\mathcal{H}_K$ in terms of the spectral data of $T_K$. Restricting to indices with $\lambda_i>0$,
        \begin{equation}\label{2605:eq:rkhs-mercer}
          \mathcal{H}_K = \Bigl\{ f=\sum_{i:\lambda_i>0}a_i e_i : \|f\|_{\mathcal{H}_K}^2 := \sum_{i:\lambda_i>0}\frac{a_i^2}{\lambda_i}<\infty \Bigr\},
          \end{equation}
          with inner product $\langle f,g\rangle_{\mathcal{H}_K}=\sum_i a_i b_i/\lambda_i$. Equivalently, $\mathcal{H}_K$ is the image of $L^2(\mathcal{X},\mu)$ under the square root $T_K^{1/2}$, with norm $\|T_K^{1/2}g\|_{\mathcal{H}_K}=\|g\|_{L^2(\mathcal{X},\mu)}$ (modulo $\ker T_K$); equivalently, $T_K^{1/2}:L^2(\mathcal{X},\mu)\to\mathcal{H}_K$ is a partial isometry. Two consequences will be used repeatedly below:
          \begin{itemize}
            \item The eigenfunctions $\{\sqrt{\lambda_i} e_i\}_{i:\lambda_i>0}$ form an orthonormal basis of $\mathcal{H}_K$, while $\{e_i\}$ form an orthonormal basis of the closure of $\mathrm{range}(T_K)$ in $L^2(\mathcal{X},\mu)$.
              \item Smoothness in $\mathcal{H}_K$ corresponds to spectral concentration: $f\in\mathcal{H}_K$ iff its $L^2$ coefficients $a_i=\langle f,e_i\rangle_{L^2}$ decay fast enough that $\sum_i a_i^2/\lambda_i<\infty$. In particular, functions in the top-$k$ eigenspace of $T_K$ are exactly the ``low-degree'' or ``low-frequency'' functions used throughout the main text.
              \end{itemize}
              This dictionary between the kernel $K$, the integral operator $T_K$ and the RKHS $\mathcal{H}_K$ is what allows us, in the rest of this appendix, to translate between function-space statements (norms, projections, low-degree truncation in $\mathcal{H}_K$) and operator-spectral statements (eigenvalues and eigenfunctions of $T_K$).
              
\subsection{Representer Property}

We begin by establishing a fundamental property: the optimal features lie in the finite-dimensional span of training features, which enables efficient computation via the kernel trick.

\begin{lemma}[Representer property of LoFi features]\label{2605:lem:representer-lofi}
    Let $\{\hat{\bm v}_j^{(\ell)}\}_{j=1}^{k_\ell}$ denote any minimizers of the empirical objective in \eqref{2605:eq:lofi-rkhs-variational}. Then the corresponding weight vectors $\{\hat{\bm v}_j^{(\ell)}\}_{j=1}^{k_\ell}$ lie in the span of the previous-layer features evaluated at the training inputs, i.e.\
    \begin{equation}\label{2605:eq:lofi-span}
        \hat{\bm v}_j^{(\ell)} \in \mathrm{span}\bigl\{\bz_{\ell-1}(\bm x_1),\dots,\bz_{\ell-1}(\bm x_n)\bigr\}, \qquad j=1,\dots,k_\ell.
    \end{equation}
\end{lemma}
\begin{proof}[Proof of \Cref{2605:lem:representer-lofi}]
Fix any index $j\in\{1,\dots,k_\ell\}$. The empirical objective in~\eqref{2605:eq:lofi-rkhs-variational} depends on $\hat{\bm v}_j^{(\ell)}$ only through the inner products $\langle \hat{\bm v}_j^{(\ell)}, \bz_{\ell-1}(\bm x_\mu)\rangle$, $\mu=1,\dots,n$, and is optimized subject to the norm constraint $\|\hat{\bm v}_j^{(\ell)}\|_2 = 1$ (together with the deflation/orthogonality constraints to previously selected directions, which are also expressible solely through such inner products). Decompose $\hat{\bm v}_j^{(\ell)} = \bm v_j^\parallel + \bm v_j^\perp$, where $\bm v_j^\parallel$ is the orthogonal projection onto the finite-dimensional subspace $\mathcal{S} := \mathrm{span}\{\bz_{\ell-1}(\bm x_\mu)\}_{\mu=1}^n \subset \mathbb{R}^{p_{\ell-1}}$ and $\bm v_j^\perp \in \mathcal{S}^\perp$. By orthogonality, $\langle \bm v_j^\perp, \bz_{\ell-1}(\bm x_\mu)\rangle = 0$ for every training input, so $\langle \hat{\bm v}_j^{(\ell)}, \bz_{\ell-1}(\bm x_\mu)\rangle = \langle \bm v_j^\parallel, \bz_{\ell-1}(\bm x_\mu)\rangle$, i.e.\ the data-dependent objective depends only on $\bm v_j^\parallel$. Pythagoras gives $1 = \|\hat{\bm v}_j^{(\ell)}\|_2^2 = \|\bm v_j^\parallel\|_2^2 + \|\bm v_j^\perp\|_2^2$, so any nonzero $\bm v_j^\perp$ forces $\|\bm v_j^\parallel\|_2 < 1$. The rescaled vector $\tilde{\bm v}_j := \bm v_j^\parallel/\|\bm v_j^\parallel\|_2 \in \mathcal{S}$ is then feasible (it has unit norm and inherits the orthogonality constraints, since these involve only inner products with vectors in $\mathcal{S}$) and yields a strictly larger objective by a factor $\|\bm v_j^\parallel\|_2^{-2}>1$, contradicting optimality of $\hat{\bm v}_j^{(\ell)}$. Hence every minimizer satisfies $\bm v_j^\perp = 0$, i.e.\ $\hat{\bm v}_j^{(\ell)} \in \mathcal{S}$, which is~\eqref{2605:eq:lofi-span}.
\end{proof}

\subsection{RKHS-Euclidean Equivalence}\label{2605:app:kernel-equiv}

We now establish a fundamental result that underlies all subsequent proofs in this section. By Lemma~\ref{2605:lem:representer-lofi}, we know features lie in the span of training features. This allows us to identify RKHS constraints with Euclidean constraints on weight vectors.

\begin{lemma}[RKHS-Euclidean Norm Equivalence]\label{2605:lem:rkhs-euclidean}
Let $K_{\ell-1}(\bm x,\bm x')=\langle\bz_{\ell-1}(\bm x),\bz_{\ell-1}(\bm x')\rangle$ and let $\mathcal H_{\ell-1}$ be its induced RKHS. Then the RKHS norm of a linear feature coincides with the Euclidean norm of its weight vector: for any $\bm u\in\mathbb R^{p_{\ell-1}}$, the linear feature
\[
    \varphi_{\bm u}({\bm x}) := \langle \bm u, \bz_{\ell-1}(\bm x)\rangle.
\]
satisfies
\[
    \|\varphi_{\bm u}\|_{\mathcal H_{\ell-1}} = \|\bm u\|_2.
\]
Consequently, the RKHS unit ball $\{\varphi:\|\varphi\|_{\mathcal H_{\ell-1}}\leq 1\}$ and the Euclidean unit sphere $\{\bm u:\|\bm u\|_2=1\}$ are in bijection through $\bm u\mapsto\varphi_{\bm u}$, establishing a primal--dual equivalence between RKHS constraints and Euclidean constraints.
\end{lemma}

\subsection{Proof of Theorem~\ref{2605:thm:variational-rkhs}}

We prove each part of Theorem~\ref{2605:thm:variational-rkhs} in turn, using Lemma~\ref{2605:lem:rkhs-euclidean} to convert between the RKHS and Euclidean formulations.

\paragraph{Proof of part (i).}

By Lemma~\ref{2605:lem:rkhs-euclidean}, the objective equals
\[
    \left|\widehat{\mathbb E}_n\bigl[y \psi(\bm x)\bigr]\right|
    =
    \left|\frac{1}{n}\sum_{\mu=1}^n y_\mu\langle \bm u, \bz_{\ell-1}(\bm x_\mu)\rangle\right|
    =
    \left|\left\langle \bm u,  \hat{\bm u}^\ell\right\rangle\right|,
\]
where $\hat{\bm u}^\ell = \tfrac{1}{n}\sum_{\mu=1}^n y_\mu \bm z_{\ell-1}(\bm x_\mu)$ is the empirical first-order moment computed in Algorithm~\ref{2605:alg:neural-lofi}. By the Cauchy--Schwarz inequality this is maximized at $\bm u = \hat{\bm u}^\ell/\|\hat{\bm u}^\ell\|_2$, giving $\psi^\ell(\bm x)=\langle \hat{\bm u}^\ell, \bm z_{\ell-1}(\bm x)\rangle$ as the unique maximizer (up to sign). \qed

\paragraph{Proof of part (ii).}

By Lemma~\ref{2605:lem:rkhs-euclidean}, the second-order objective equals
\[
    \left|\widehat{\mathbb E}_n\bigl[y \varphi(\bm x)^2\bigr]\right|
    =
    \left|\frac{1}{n}\sum_{\mu=1}^n y_\mu\langle \bm u,\bz_{\ell-1}(\bm x_\mu)\rangle^2\right|
    =
    \left|\bm u^\top \widehat{\bm C}^{(\ell)} \bm u\right|,
\]
where $\widehat{\bm C}^{(\ell)} = \tfrac{1}{n}\sum_{\mu=1}^n y_\mu \bz_{\ell-1}(\bm x_\mu)\bz_{\ell-1}(\bm x_\mu)^\top$ is the empirical second-order moment operator. By the Courant--Fischer minimax theorem, the unit-norm vector maximizing $|\bm u^\top \widehat{\bm C}^{(\ell)}\bm u|$ is the leading eigenvector $\hat{\bm v}^{(\ell)}_1$ of $\widehat{\bm C}^{(\ell)}$ (in absolute eigenvalue). Successive orthogonal maximizers are the subsequent eigenvectors $\hat{\bm v}^{(\ell)}_2,\dots,\hat{\bm v}^{(\ell)}_{k_\ell}$. Via $\bm u\mapsto\varphi_{\bm u}$, these correspond exactly to $\varphi_1,\dots,\varphi_{k_\ell}$ satisfying the stated variational recursion. \qed

\paragraph{Dual (kernel-space) form ---}

By Lemma~\ref{2605:lem:representer-lofi}, any optimal feature $\varphi_j$ lies in the span of the $n$ training feature vectors:
\begin{equation}
\varphi_j(\bm x) = \sum_{\mu=1}^n \alpha_{j,\mu} K_{\ell-1}(\bm x, \bm x_\mu).
\end{equation}
Define $\bm\alpha_j = (\alpha_{j,1}, \dots, \alpha_{j,n})^\top$ as the coefficient vector for feature $j$. By Lemma~\ref{2605:lem:rkhs-euclidean}, the RKHS norm of this feature is
\begin{equation}
\|\varphi_j\|^2_{\mathcal H_{\ell-1}} = \bm\alpha_j^\top G_{\ell-1} \bm\alpha_j,
\end{equation}
where $G_{\ell-1} = Z_{\ell-1}Z_{\ell-1}^\top \in \mathbb R^{n \times n}$ is the Gram matrix with entries $G_{\ell-1}(\mu, \nu) = K_{\ell-1}(\bm x_\mu, \bm x_\nu)$.

The empirical objective from Theorem~\ref{2605:thm:variational-rkhs}~(ii) becomes
\begin{equation}
\max_{\bm\alpha:\|\varphi\|_{\mathcal H_{\ell-1}}=1} \left|\frac{1}{n}\sum_{\mu=1}^n y_\mu \varphi(\bm x_\mu)^2\right| = \max_{\bm\alpha:\bm\alpha^\top G_{\ell-1}\bm\alpha=1} \left|\bm\alpha^\top \widehat{\bm C}^{(\ell)} \bm\alpha\right|,
\end{equation}
where $\widehat{\bm C}^{(\ell)} = \frac{1}{n}\sum_{\mu=1}^n y_\mu K_{\ell-1}(\bm x_\mu, \cdot) \otimes K_{\ell-1}(\bm x_\mu, \cdot)$ is the label-weighted kernel outer-product operator.

\begin{proposition}[Generalized Eigenvector Problem]\label{2605:prop:generalized-eigenvector}
The optimal dual coefficients $\{\hat{\bm\alpha}_j\}_{j=1}^{k_\ell}$ that sequentially maximize the above constrained objective satisfy the generalized eigenvector problem
\begin{equation}
G_{\ell-1}^{1/2} Y G_{\ell-1}^{1/2} \bm\alpha_j = \lambda_j \bm\alpha_j, \quad j=1,\dots,k_\ell,
\label{2605:eq:gen-eigenproblem}
\end{equation}
where $Y = \mathrm{diag}(y_1, \dots, y_n)$ is the label matrix, and $\lambda_j$ are the generalized eigenvalues ordered by magnitude. The solutions $\{\hat{\bm\alpha}_j\}_{j=1}^{k_\ell}$ recover the dual coefficients of the second-order features, enabling a kernel implementation of Neural LoFi without explicit access to the feature vectors.
\end{proposition}

\paragraph{Proof of part (iii).}
At every layer, assume that the limiting eigenvalue (or
eigenvalue cluster) being selected is separated from the rest of the spectrum. For a single
feature this means a positive eigengap; for a cluster the statements below hold for the
spectral projector, with an arbitrary orthonormal basis chosen inside the limiting
eigenspace.

\emph{Induction statement.} Let $\rho_{\bm x}$ be the data distribution. After layer $r$, let
\[
    \bm g_r^{(p)}(\bm x)
    =
    \bigl(\psi_r^{(p)}(\bm x),\varphi_{r,1}^{(p)}(\bm x),\dots,
    \varphi_{r,k_r}^{(p)}(\bm x)\bigr)
\]
denote the projected features selected by Neural LoFi at finite width, and let
$K_r^{(p)}$ be the kernel produced by the following random lift. The induction claim
$\mathsf I_r$ is that there exist deterministic functions $\bm g_r^\infty$ and a
deterministic kernel $K_r^\infty$ such that
\begin{equation}
    \|\bm g_r^{(p)}-\bm g_r^\infty\|_{L^2(\rho_{\bm x})}\xrightarrow[p\to\infty]{P}0
    \label{2605:eq:feature-induction}
\end{equation}
and, for the training inputs,
\begin{equation}
    \|\bm G_r^{(p)}-\bm G_r^\infty\|_{\mathrm{op}}\xrightarrow[p\to\infty]{P}0,\qquad
    \sum_{\mu=1}^n
    \|K_r^{(p)}(\cdot,\bm x_\mu)-K_r^\infty(\cdot,\bm x_\mu)\|_{L^2(\rho_{\bm x})}^2
    \xrightarrow[p\to\infty]{P}0,
    \label{2605:eq:kernel-section-induction}
\end{equation}
where $(\bm G_r^{(p)})_{\mu\nu}=K_r^{(p)}(\bm x_\mu,\bm x_\nu)$.
The base case is deterministic: $\bm g_0(\bm x)=\bm x$ and $K_0(\bm x,\bm x')=
\langle\bm x,\bm x'\rangle$.

We first record the kernel-propagation step used in the induction. Suppose that, for some
layer $r$, the projected features already satisfy
$\bm g_r^{(p)}\to\bm g_r^\infty$ in $L^2(\rho_{\bm x})$. For a deterministic feature map
$\bm g$, write
\[
    \mathcal K_r[\bm g](\bm x,\bm x')
    =
    \E_{\bm a\sim\pi_r}
    \bigl[
        \sigma(\bm a^\top \bm g(\bm x))
        \sigma(\bm a^\top \bm g(\bm x'))
    \bigr].
\]
The limiting next-layer kernel is
$K_r^\infty=\mathcal K_r[\bm g_r^\infty]$. Since $\bm g_r^{(p)}\to \bm g_r^\infty$ in
$L^2$ and $\sigma$ is pseudo-Lipschitz with the required moment bounds, dominated
convergence gives convergence of the deterministic kernels
$\mathcal K_r[\bm g_r^{(p)}]\to \mathcal K_r[\bm g_r^\infty]$ on the fixed Gram matrix
and in the $L^2$ section norm in~\eqref{2605:eq:kernel-section-induction}. Conditional on
$\bm g_r^{(p)}$, the finite-width kernel is an average of iid random features:
\[
    K_r^{(p)}(\bm x,\bm x')
    =
    \frac1{p_r}\sum_{i=1}^{p_r}
    \sigma(\bm a_i^\top \bm g_r^{(p)}(\bm x))
    \sigma(\bm a_i^\top \bm g_r^{(p)}(\bm x')).
\]
For each fixed anchor $\bm x_\mu$,
\begin{equation}
    \E_{\bm a}
    \bigl\|
        K_r^{(p)}(\cdot,\bm x_\mu)
        -\mathcal K_r[\bm g_r^{(p)}](\cdot,\bm x_\mu)
    \bigr\|_{L^2(\rho_{\bm x})}^2
    \leq \frac{C_\mu}{p_r},
    \label{2605:eq:rf-section-lln}
\end{equation}
and the same iid law of large numbers applies to the finitely many Gram entries. Thus
\eqref{2605:eq:kernel-section-induction} follows. This is the formal sense in which, once the
features entering a layer have a deterministic limit, the next layer sees a fixed
deterministic distribution.

It remains to show that this deterministic kernel convergence propagates through the
spectral filtering step. Assume $\mathsf I_{\ell-1}$ and abbreviate
$K_p=K_{\ell-1}^{(p)}$, $K_\infty=K_{\ell-1}^\infty$, and
$\bm G_p,\bm G_\infty$ for their Gram matrices. By~\eqref{2605:eq:kernel-section-induction},
\begin{equation}
    \|\bm G_p-\bm G_\infty\|_{\mathrm{op}}\xrightarrow[p\to\infty]{P}0,\qquad
    \sum_{\mu=1}^n
    \|K_p(\cdot,\bm x_\mu)-K_\infty(\cdot,\bm x_\mu)\|_{L^2(\rho_{\bm x})}^2
    \xrightarrow[p\to\infty]{P}0.
    \label{2605:eq:gram-and-section-conv}
\end{equation}
The explicit $O_P(n/\sqrt p)$ Gram-matrix rate in the non-adaptive, fixed-$\bm g$
case is the usual random-feature concentration bound; for the section convergence,
\eqref{2605:eq:rf-section-lln} gives the sharper Hilbert-space law of large numbers needed for
out-of-sample features.

By Lemma~\ref{2605:lem:representer-lofi}, every selected feature has the form
\[
    \varphi(\bm x)=\sum_{\mu=1}^n \alpha_\mu K_p(\bm x,\bm x_\mu).
\]
Its RKHS norm and empirical second-order objective are
\[
    \|\varphi\|_{\cH_p}^2=\balpha^\top \bm G_p\balpha,\qquad
    \widehat{\E}_n[y\varphi(\bm x)^2]
    =
    \frac1n\balpha^\top \bm G_p Y \bm G_p\balpha,
    \qquad
    Y=\operatorname{diag}(y_1,\dots,y_n).
\]
Therefore, on the range of $\bm G_p$, the generalized eigenproblem is equivalently the
symmetric eigenproblem
\begin{equation}
    \bm B_p \bm\beta=\lambda \bm\beta,\qquad
    \bm B_p
    =
    \frac1n \bm G_p^{1/2}Y\bm G_p^{1/2},\qquad
    \bm\beta=\bm G_p^{1/2}\balpha.
    \label{2605:eq:symmetric-dual-eigenproblem}
\end{equation}
The square-root map is continuous on positive semidefinite matrices, hence
$\|\bm B_p-\bm B_\infty\|_{\mathrm{op}}\to 0$ in probability. If the $k$th limiting eigenvalue is
isolated with gap $\Delta_k>0$, the Davis--Kahan $\sin\Theta$ theorem
\cite{davis1970rotation} implies convergence of the corresponding projectors at
rate $O_P(\|\bm B_p-\bm B_\infty\|_{\mathrm{op}}/\Delta_k)$. For a simple eigenvalue, after fixing
the sign,
\begin{equation}
    \bm\beta_{p,k}\xrightarrow[p\to\infty]{P}\bm\beta_{\infty,k}.
    \label{2605:eq:beta-conv}
\end{equation}


We next convert the convergence of $\bm\beta_{p,k}$ into convergence of the actual feature functions.
Let
\[
    \bm k_p(\bm x)=\bigl(K_p(\bm x,\bm x_1),\dots,K_p(\bm x,\bm x_n)\bigr)^\top,
    \qquad
    \balpha_{p,k}=\bm G_p^{\dagger/2}\bm\beta_{p,k},
\]
where $\dagger$ denotes the Moore--Penrose inverse on the stable range of the limiting
Gram matrix. If $\bm G_\infty$ is nonsingular, this is the ordinary inverse square root. Continuity of the pseudo-inverse operator and~\eqref{2605:eq:beta-conv} give
$\balpha_{p,k}\to\balpha_{\infty,k}$ in probability. Hence, with
\[
    \widehat\varphi_{p,k}(\bm x)=\bm k_p(\bm x)^\top\balpha_{p,k},\qquad
    \phi^\infty_k(\bm x)=\bm k_\infty(\bm x)^\top\balpha_{\infty,k},
\]
we have
\begin{equation}
\begin{aligned}
    \|\widehat\varphi_{p,k}-\phi^\infty_k\|_{L^2(\rho_{\bm x})}
    &\leq
    \|\bm k_p-\bm k_\infty\|_{L^2(\rho_{\bm x};\R^n)} 
    \|\balpha_{p,k}\|_2 \\
    &\quad+
    \|\bm k_\infty\|_{L^2(\rho_{\bm x};\R^n)} 
    \|\balpha_{p,k}-\balpha_{\infty,k}\|_2
    \xrightarrow[p\to\infty]{P}0.
\end{aligned}
    \label{2605:eq:feature-function-conv}
\end{equation}
The linear feature satisfies the same conclusion directly, since
\[
    \psi_p^\ell(\bm x)
    =
    \frac1n\sum_{\mu=1}^n y_\mu K_p(\bm x,\bm x_\mu)
    \to
    \frac1n\sum_{\mu=1}^n y_\mu K_\infty(\bm x,\bm x_\mu)
    =
    \psi_\infty^\ell(\bm x)
\]
in $L^2(\rho_{\bm x})$. Thus the whole projected vector
$\bm g_\ell^{(p)}=(\psi_p^\ell,\widehat\varphi_{p,1},\dots,\widehat\varphi_{p,k_\ell})$
converges to the deterministic vector $\bm g_\ell^\infty$ in $L^2$. Applying the first
part of the induction to this deterministic limit gives convergence of the next random
feature kernel $K_\ell^{(p)}$ to
\[
    K_\ell^\infty(\bm x,\bm x')
    =
    \E_{\bm a\sim\pi_\ell}
    \left[
        \sigma(\bm a^\top\bm g_\ell^\infty(\bm x))
        \sigma(\bm a^\top\bm g_\ell^\infty(\bm x'))
    \right],
\]
which proves $\mathsf I_\ell$. By induction, all finite collections of learned features and
the induced kernels converge layer by layer to deterministic infinite-width limits.


 \qed
